Practical physics education requires environments that support analysis, review of decisions, and preservation of results. This technical essay presents an architecture based on Git, development containers, notebooks, and automation. Documentary analysis compares Codespaces, Binder, JupyterHub, and Colab in terms of administration, persistence, and operational constraints, contextualized in Latin American teaching materials. Results distinguish repository availability, verifiable execution, and session evaluation as separate objects. A proposed protocol separates instructor preparation, student onboarding, resource consumption, and scientific interpretation, supported by recording templates. The analysis explains how infrastructure can be selected according to the task and available support, and how reproducible practices can transfer to other activities. The editorial repository provides an executable example of the same procedure.
